\chapter{OPERATOR KOMPAK}\label{chap_cpt_ops}

\section{Definisi dan Sifat-Sifat Dasar}

Kita akan melanjutkan pembahasan dari akhir Bab 5 dengan mempelajari kelas-kelas operator.
Konteks kita akan sedikit lebih umum; dalam hal operator kompak, kita akan
meninjau operator pada ruang Banach.  Pertama, mari kita ingat kembali beberapa definisi dan
fakta dasar dari kalkulus lanjut (atau mata kuliah pengantar analisis).

\begin{defn} Suatu subhimpunan $A$ dari ruang metrik $M$ disebut
 \index{total!keterbatasan}%
 \index{terbatas!total}%
\df{terbatas total} jika untuk setiap $\epsilon > 0$ terdapat suatu subhimpunan hingga $F$ dari $A$
sedemikian sehingga untuk setiap $a \in A$ terdapat titik $x \in F$ yang memenuhi $d(x,a) < \epsilon$.
Untuk ruang linear bernorma, definisi ini dapat dirumuskan kembali: Suatu subhimpunan $A$ dari ruang linear
bernorma $V$ disebut \df{terbatas total} jika untuk setiap bola terbuka $B$ di sekitar nol dalam $V$ terdapat
suatu subhimpunan hingga $F$ dari $A$ sedemikian sehingga $A \subseteq F + B$.  Hal ini memiliki parafrasa yang kurang lebih baku:
Suatu ruang bersifat terbatas total jika dapat diawasi oleh sejumlah hingga polisi dengan jarak pandang
sependek apa pun.  (Sebagian penulis menyebut sifat ini \emph{prakompak}.)
\end{defn}

\begin{prop} Suatu ruang metrik bersifat kompak jika dan hanya jika ruang itu lengkap dan terbatas total.
\end{prop}

\begin{prop} Suatu ruang metrik bersifat kompak jika dan hanya jika ruang itu lengkap dan terbatas total.
\end{prop}

\begin{defn} Suatu subhimpunan $A$ dari ruang topologis $X$ disebut
 \index{kompak!relatif}%
 \index{relatif kompak}%
\df{kompak relatif} jika tutupannya kompak.
\end{defn}

\begin{prop} Dalam ruang metrik, setiap himpunan kompak relatif adalah terbatas total.
\end{prop}

Dalam ruang metrik lengkap, kebalikannya berlaku.

\begin{prop} Dalam ruang metrik lengkap, setiap himpunan terbatas total adalah kompak relatif.
\end{prop}

\begin{prop} Dalam ruang metrik, setiap himpunan terbatas total adalah separabel.
\end{prop}

\begin{thm}[Teorema Arzel\`a-Ascoli]\label{AAThm} Misalkan $X$ ruang Hausdorff kompak. Suatu subhimpunan $\fml F$
dari $\fml C(X)$ terbatas total apabila subhimpunan itu terbatas (secara seragam) dan ekuikontinu.
\end{thm}

Berikutnya kita akan membandingkan topologi \emph{lemah} dan \emph{kuat} pada ruang linear bernorma.
Istilah
 \index{lemah!\emph{vs.} topologi kuat pada ruang linear bernorma}%
 \index{topologi!lemah!\emph{vs.} kuat pada ruang linear bernorma}%
 \index{kuat!\emph{vs.} topologi lemah pada ruang linear bernorma}%
 \index{konvensi!pada ruang linear bernorma \emph{topologi kuat} merujuk pada topologi norma}%
\emph{topologi lemah} akan merujuk pada topologi lemah biasa yang ditentukan oleh fungsional linear terbatas
pada ruang tersebut sebagaimana didefinisikan dalam~\ref{C067434}, sedangkan \emph{topologi kuat} akan merujuk pada topologi norma pada ruang tersebut.
Namun, perlu diingat bahwa kata ``kuat'' dalam konteks ruang linear bernorma, jika ditilik secara ketat,
sesungguhnya mubazir.  Kata itu digunakan semata-mata sebagai penegasan.  Misalnya, banyak penulis akan menyatakan proposisi berikut sebagai:
\emph{Setiap barisan yang konvergen secara lemah dalam ruang linear bernorma adalah terbatas.}

\begin{prop} Setiap barisan yang konvergen secara lemah dalam ruang linear bernorma adalah terbatas secara kuat.
\end{prop}

\begin{proof}[\emph{Petunjuk untuk bukti}] Gunakan \emph{prinsip keterbatasan seragam}~\ref{thm_PUB}. \ns
\end{proof}

\begin{defn} Kita mengatakan bahwa suatu pemetaan linear $T\colon V \sto W$ bersifat
 \index{lemah!kekontinuan}%
 \index{kontinu!secara lemah}%
 \index{linear!pemetaan!kekontinuan lemah suatu}%
\df{kontinu secara lemah} jika pemetaan itu membawa jaring yang konvergen secara lemah ke jaring yang konvergen secara lemah.
\end{defn}

\begin{prop} Setiap pemetaan linear yang kontinu (secara kuat) antara ruang-ruang linear bernorma bersifat kontinu secara lemah.
\end{prop}

\begin{defn} Suatu pemetaan linear $T\colon V \sto W$ antara ruang-ruang linear bernorma disebut
 \index{kompak!pemetaan linear}%
 \index{linear!pemetaan!kompak}%
 \index{operator!kompak}%
\df{kompak} jika pemetaan itu membawa himpunan terbatas ke himpunan kompak relatif. Secara ekuivalen, $T$ bersifat kompak jika pemetaan itu
memetakan bola satuan tertutup dalam $V$ ke suatu himpunan kompak relatif dalam~$W$.   Keluarga semua pemetaan linear kompak
dari $V$ ke $W$ dinotasikan
 \index{K@$\ofml K(V,W)$!keluarga pemetaan linear kompak}%
 \index{K@$\ofml K(V)$!keluarga operator kompak}%
dengan~$\ofml K(V,W)$; dan kita tulis $\ofml K(V)$ untuk~$\ofml K(V,V)$.
\end{defn}

\begin{prop} Setiap pemetaan linear kompak bersifat terbatas.
\end{prop}

\begin{prop} Suatu pemetaan linear $T\colon V \sto W$ antara ruang-ruang linear bernorma bersifat kompak jika dan hanya jika pemetaan itu membawa
setiap barisan terbatas dalam $V$ ke suatu barisan dalam $W$ yang memiliki subbarisan konvergen.
\end{prop}

\begin{exam} Setiap operator berperingkat hingga pada ruang Banach bersifat kompak.
\end{exam}

\begin{exam} Setiap fungsional linear terbatas pada ruang Banach $B$ merupakan pemetaan linear kompak dari $B$ ke~$\K$.
\end{exam}

\begin{exam} Misalkan $C$ ruang Banach $\fml C([0,1])$ dan $k \in \fml C([0,1] \times [0,1])$.  Untuk setiap $f \in C$
 \index{integral!operator}%
 \index{operator!integral}%
definisikan fungsi $Kf$ pada $[0,1]$ dengan
   \[ Kf(x) = \int_0^1 k(x,y)f(y)\,dy\,. \]
Maka operator integral $K$ merupakan operator kompak pada~$C$.
\end{exam}

\begin{proof}[\emph{Petunjuk untuk bukti}] Gunakan \emph{teorema Arzel\`a-Ascoli}~\ref{AAThm}.  \ns
\end{proof}

\begin{exam}\label{X_square_int_kernel} Misalkan $H$ ruang Hilbert $\lfs2{[0,1]}$ yang terdiri atas (kelas-kelas ekuivalensi) fungsi yang
terintegralkan kuadrat pada $[0,1]$ terhadap ukuran Lebesgue dan
 \index{integral!operator}%
 \index{operator!integral}%
$k$ suatu fungsi yang terintegralkan kuadrat
pada $[0,1] \times [0,1]$.  Untuk setiap $f \in H$ definisikan fungsi $Kf$ pada $[0,1]$ dengan
   \[ Kf(x) = \int_0^1 k(x,y)f(y)\,dy\,. \]
Maka operator integral $K$ merupakan operator kompak pada~$H$.
\end{exam}

\begin{prop} Suatu operator pada ruang Hilbert bersifat kompak jika dan hanya jika operator itu memetakan bola satuan tertutup dalam ruang tersebut
ke suatu himpunan kompak.
\end{prop}

\begin{cor}\label{cor_id_op_not_cpt} Operator identitas pada ruang Hilbert berdimensi tak hingga tidak kompak.
(Lihat contoh~\ref{X_l2ball_not_cpt}.)
\end{cor}

\begin{exam} Jika $B$ ruang Banach, keluarga $\ofml B(B)$ dari operator pada $B$ merupakan
 \index{Banach!aljabar!$\ofml B(B)$ sebagai suatu}%
 \index{B@$\ofml B(V)$!sebagai aljabar Banach beridentitas}%
aljabar Banach beridentitas dan
 \index{K@$\ofml K(V)$!sebagai ideal tertutup sejati}%
keluarga $\ofml K(B)$ dari operator kompak pada $B$ merupakan ideal tertutup dalam~$\ofml B(B)$.
Jika $B$ berdimensi tak hingga, ideal $\ofml K(B)$ adalah sejati.
\end{exam}

\begin{prop} Misalkan $T\colon B \sto C$ suatu pemetaan linear terbatas antara ruang-ruang Banach. Maka $T$ bersifat kompak
jika dan hanya jika $T^*$ bersifat kompak.
\end{prop}
